% paper/sections/appendix_fvm_note.tex
%
% §8 補足：FVM 基準との整合（調和平均の物理的根拠）
%
%   本節は参照用補足である．本稿の PPE ソルバーは CCD product-rule 展開を
%   採用しており，FVM は使用しない．調和平均の物理的根拠を記録する目的で残す．
%   ※ Rhie--Chow 補正済み発散の定義・チェッカーボード解説・零モード対策・ピン点選択は
%     §8.4（08f_ppe_bc.tex）および付録 F（appendix_legacy_methods.tex）に統合済み（CHK-087）．
% ═══════════════════════════════════════════════════════════════════════════════
\subsection{【補足】FVM 基準との整合：調和平均の物理的根拠}
\label{sec:fvm_face_coeff}

PPE 演算子 $\bnabla\cdot(a\bnabla p)$（$a \equiv 1/\rho$）を FVM で離散化すると，
フェイス面係数として\textbf{調和平均}が自然に現れる（積分論理は付録~\ref{app:fvm_face_coeff}）：
\begin{equation}
  a_f = \frac{2}{\rho_L + \rho_R}
  \label{eq:af_exact}
\end{equation}
よくある実装ミス：$\rho$ の調和平均の逆数 $(\rho_L+\rho_R)/(2\rho_L\rho_R) \neq 2/(\rho_L+\rho_R)$．
\textbf{本稿の実装}は FVM ではなく CCD product-rule 展開（式~\eqref{eq:varrho_product_rule}，§\ref{sec:ccd_poisson_matrix}）であり，$\Ord{h^6}$ を達成する．
式~\eqref{eq:af_exact} は GFM 補正項の調和平均密度逆数（式~\eqref{eq:gfm_rhs_correction}）と一致する（§\ref{sec:gfm}）．

Rhie--Chow 補正済み発散 $\bnabla_h^{\mathrm{RC}}\!\cdot\bu^*$ の完全な定義，
チェッカーボード安定性との関係，および旧 Balanced-Force 定式化は
付録~\ref{app:rhie_chow_full} に収録されている．
本稿は DCCD 散逸（§\ref{sec:dccd_decoupling}）+ GFM（§\ref{sec:gfm}）により
Rhie--Chow 補正を代替する（式~\eqref{eq:gfm_ccd_div}）．
